\begin{enumerate}
  \item On modélise le dispositif du jeu par un solide ( $S$ ) de centre d'inertie $G$ et de masse $m$ se déplaçant sur un rail rectiligne $A B$ de longueur $L$ incliné d'un angle $\alpha$ par rapport à l'horizontal (Figure 1). À l'instant $t_{0}=0$, le solide $(S)$ est lancé avec une vitesse initiale $\vec{V}_{0}$ de la position
\end{enumerate}

\begin{figure}[h]
\begin{center}
  \includegraphics[alt={},max width=\textwidth]{53bec4bf-ae8a-49fa-b6a9-4c685324fa09-5_341_533_1631_1361}
\captionsetup{labelformat=empty}
\caption{Figure 1}
\end{center}
\end{figure}

$A$.\\
On étudie le mouvement de ( $S$ ) dans un repère ( $O, \vec{i}, \vec{j}$ ) lié à la Terre supposé galiléen.\\
L'abscisse de $G$ à $t_{0}=0$ est $x_{G}=x_{0}=0$.\\
Données : $m=1 \mathrm{~kg} ; \alpha=25^{\circ} ; g=10 \mathrm{~m} . \mathrm{s}^{-2} ; L=A B=3 \mathrm{~m}$\\
Au cours de son mouvement, le solide subit des frottements équivalents à une force $\bar{f}$ constante colinéaire au vecteur vitesse et de sens opposé.\\
1.1. En appliquant la deuxième loi de Newton, montrer que l'expression de l'accélération $a_{G}$ de $G$ s'écrit : $a_{G}=-\left(g \cdot \sin \alpha+\frac{f}{m}\right)$.\\
1.2. La figure 2 donne le diagramme de vitesse de $G$.

\begin{figure}[h]
\begin{center}
  \includegraphics[alt={},max width=\textwidth]{53bec4bf-ae8a-49fa-b6a9-4c685324fa09-5_538_543_2051_1356}
\captionsetup{labelformat=empty}
\caption{Figure 2}
\end{center}
\end{figure}

1.2.1. Écrire l'expression numérique de l'équation de la vitesse de $G$.\\
1.2.2. Calculer l'intensité $f$ de la force de frottement.\\
1.2.3. Déterminer l'instant $t_{1}$ où le sens de mouvement de $(S)$ change.\\
1.2.4. Calculer la distance $d_{1}$ parcourue par le solide durant la phase de montée.\\
1.2.5. Déterminer la vitesse minimale $v_{0, \text { min }}$ avec laquelle le solide $(S)$ doit être lancé pour réaliser le défi.\\

